\begin{codeblock}
namespace std {
  template<class T>
    requires move_constructible<T> && is_object_v<T>
  class @\libglobal{movable_box}@ {
  public:
    using @\libmember{value_type}{movable_box}@ = T;

    // \ref{movable.box.ctor}, constructors
    constexpr movable_box() noexcept(is_nothrow_default_constructible_v<T>)
      requires default_initializable<T>;
    constexpr explicit movable_box(const T& t)
      requires copy_constructible<T>;
    constexpr explicit movable_box(T&& t);
    template<class... Args>
      requires is_constructible_v<T, Args...>
    constexpr explicit movable_box(in_place_t, Args&&... args);

    constexpr movable_box(const movable_box& rhs)
      requires copy_constructible<T>;
    constexpr movable_box(movable_box&& rhs) noexcept(
        is_nothrow_move_constructible_v<T>);

    // \ref{movable.box.dtor}, destructor
    constexpr ~movable_box();

    // \ref{movable.box.assign}, assignment
    constexpr movable_box& operator=(const movable_box& rhs) noexcept(@\seebelow@)
      requires copy_constructible<T>;
    constexpr movable_box& operator=(movable_box&& rhs) noexcept(@\seebelow@);

    // \ref{movable.box.obs}, observers
    constexpr bool has_value() const noexcept;
    constexpr explicit operator bool() const noexcept;

    constexpr T& operator*() & noexcept;
    constexpr const T& operator*() const& noexcept;
    constexpr T&& operator*() && noexcept;
    constexpr const T&& operator*() const&& noexcept;

    constexpr T* operator->() noexcept;
    constexpr const T* operator->() const noexcept;
  };
}
\end{codeblock}

\rSec3[movable.box.ctor]{Constructors}

\indexlibraryctor{movable_box}%
\begin{itemdecl}
constexpr movable_box() noexcept(is_nothrow_default_constructible_v<T>)
  requires default_initializable<T>;
\end{itemdecl}

\begin{itemdescr}
\pnum
\effects
Equivalent to \tcode{movable_box(in_place)}.
\end{itemdescr}

\indexlibraryctor{movable_box}%
\begin{itemdecl}
constexpr explicit movable_box(const T& t)
  requires copy_constructible<T>;
\end{itemdecl}

\begin{itemdescr}
\pnum
\effects
Direct-non-list-initializes the contained value with \tcode{t}.
\end{itemdescr}

\indexlibraryctor{movable_box}%
\begin{itemdecl}
constexpr explicit movable_box(T&& t);
\end{itemdecl}

\begin{itemdescr}
\pnum
\effects
Direct-non-list-initializes the contained value with \tcode{std::move(t)}.
\end{itemdescr}

\indexlibraryctor{movable_box}%
\begin{itemdecl}
template<class... Args> constexpr explicit movable_box(in_place_t, Args&&... args);
\end{itemdecl}

\begin{itemdescr}
\pnum
\constraints
\tcode{is_constructible_v<T, Args...>} is \tcode{true}.

\pnum
\effects
Direct-non-list-initializes the contained value with \tcode{std::forward<Args>(args)...}.
\end{itemdescr}

\indexlibraryctor{movable_box}%
\begin{itemdecl}
constexpr movable_box(const movable_box& rhs)
  requires copy_constructible<T>;
\end{itemdecl}

\begin{itemdescr}
\pnum
\effects
If \tcode{rhs} contains a value, direct-non-list-initializes the contained value with \tcode{*rhs}; otherwise, \tcode{*this} contains no value.

\pnum
\remarks
This constructor is trivial if \tcode{is_trivially_copy_constructible_v<T>} is \tcode{true}.
\end{itemdescr}

\indexlibraryctor{movable_box}%
\begin{itemdecl}
constexpr movable_box(movable_box&& rhs) noexcept(is_nothrow_move_constructible_v<T>);
\end{itemdecl}

\begin{itemdescr}
\pnum
\effects
If \tcode{rhs} contains a value, direct-non-list-initializes the contained value with \tcode{std::move(*rhs)}; otherwise, \tcode{*this} contains no value. \tcode{rhs.has_value()} is unchanged.

\pnum
\remarks
This constructor is trivial if \tcode{is_trivially_move_constructible_v<T>} is \tcode{true}.
\end{itemdescr}

\rSec3[movable.box.dtor]{Destructor}

\indexlibrarydtor{movable_box}%
\begin{itemdecl}
constexpr ~movable_box();
\end{itemdecl}

\begin{itemdescr}
\pnum
\effects
Destroys the contained value, if any.

\pnum
\remarks
This destructor is trivial if \tcode{is_trivially_destructible_v<T>} is \tcode{true}.
\end{itemdescr}

\rSec3[movable.box.assign]{Assignment}

\indexlibrarymember{operator=}{movable_box}%
\begin{itemdecl}
constexpr movable_box& operator=(const movable_box& rhs) noexcept(@\seebelow@)
  requires copy_constructible<T>;
\end{itemdecl}

\begin{itemdescr}
\pnum
\effects
If \tcode{this == addressof(rhs)} is \tcode{true}, no effects. Otherwise:
\begin{itemize}
\item If \tcode{copyable<T>} is modeled: if \tcode{*this} and \tcode{rhs} each contain a value, the contained value of \tcode{*this} is copy-assigned from \tcode{*rhs}; otherwise, if \tcode{rhs} contains a value, a contained value is direct-non-list-initialized with \tcode{*rhs}; otherwise, the contained value of \tcode{*this}, if any, is destroyed.
\item Otherwise: the contained value of \tcode{*this}, if any, is destroyed; then, if \tcode{rhs} contains a value, a contained value is direct-non-list-initialized with \tcode{*rhs}. If the initialization exits via an exception, \tcode{*this} contains no value.
\end{itemize}

\pnum
\returns
\tcode{*this}.

\pnum
\remarks
The exception specification is equivalent to \tcode{is_nothrow_copy_constructible_v<T>} and, if \tcode{copyable<T>} is modeled, \tcode{is_nothrow_copy_assignable_v<T>}.

\pnum
This operator is trivial if \tcode{is_trivially_copy_constructible_v<T>}, \tcode{is_trivially_copy_assignable_v<T>}, and \tcode{is_trivially_destructible_v<T>} are all \tcode{true}.
\end{itemdescr}

\indexlibrarymember{operator=}{movable_box}%
\begin{itemdecl}
constexpr movable_box& operator=(movable_box&& rhs) noexcept(@\seebelow@);
\end{itemdecl}

\begin{itemdescr}
\pnum
\effects
If \tcode{this == addressof(rhs)} is \tcode{true}, no effects. Otherwise:
\begin{itemize}
\item If \tcode{movable<T>} is modeled: if \tcode{*this} and \tcode{rhs} each contain a value, the contained value of \tcode{*this} is move-assigned from \tcode{*rhs}; otherwise, if \tcode{rhs} contains a value, a contained value is direct-non-list-initialized with \tcode{std::move(*rhs)}; otherwise, the contained value of \tcode{*this}, if any, is destroyed. \tcode{rhs.has_value()} is unchanged.
\item Otherwise: the contained value of \tcode{*this}, if any, is destroyed; then, if \tcode{rhs} contains a value, a contained value is direct-non-list-initialized with \tcode{std::move(*rhs)}. If the initialization exits via an exception, \tcode{*this} contains no value. \tcode{rhs.has_value()} is unchanged.
\end{itemize}

\pnum
\returns
\tcode{*this}.

\pnum
\remarks
The exception specification is equivalent to \tcode{is_nothrow_move_constructible_v<T>} and, if \tcode{movable<T>} is modeled, \tcode{is_nothrow_move_assignable_v<T>}.

\pnum
This operator is trivial if \tcode{is_trivially_move_constructible_v<T>}, \tcode{is_trivially_move_assignable_v<T>}, and \tcode{is_trivially_destructible_v<T>} are all \tcode{true}.
\end{itemdescr}

\rSec3[movable.box.obs]{Observers}

\indexlibrarymember{has_value}{movable_box}%
\indexlibrarymember{operator bool}{movable_box}%
\begin{itemdecl}
constexpr bool has_value() const noexcept;
constexpr explicit operator bool() const noexcept;
\end{itemdecl}

\begin{itemdescr}
\pnum
\returns
\tcode{true} if and only if \tcode{*this} contains a value.
\end{itemdescr}

\indexlibrarymember{operator*}{movable_box}%
\begin{itemdecl}
constexpr T& operator*() & noexcept;
constexpr const T& operator*() const& noexcept;
\end{itemdecl}

\begin{itemdescr}
\pnum
\expects
\tcode{*this} contains a value.

\pnum
\returns
A reference to the contained value.
\end{itemdescr}

\indexlibrarymember{operator*}{movable_box}%
\begin{itemdecl}
constexpr T&& operator*() && noexcept;
constexpr const T&& operator*() const&& noexcept;
\end{itemdecl}

\begin{itemdescr}
\pnum
\expects
\tcode{*this} contains a value.

\pnum
\returns
\tcode{std::move(}the contained value\tcode{)}.
\end{itemdescr}

\indexlibrarymember{operator->}{movable_box}%
\begin{itemdecl}
constexpr T* operator->() noexcept;
constexpr const T* operator->() const noexcept;
\end{itemdecl}

\begin{itemdescr}
\pnum
\expects
\tcode{*this} contains a value.

\pnum
\returns
\tcode{addressof(}the contained value\tcode{)}.
\end{itemdescr}
